%%%%%%%%%%%%%%%%%%%%%%%%%%%%%%%%%%%%%%%%%%%%%%%%%%%%%%%%%%%%%%%%%%%%%%%%
%% Art. 27, VI — descrever de forma precisa, clara e suficiente, de     %%
%% maneira que um técnico no assunto possa realizá-la, fazendo remissão %%
%% aos sinais de referência constantes dos desenhos.                    %%
%%                                                                     %%
%% Esta é a maior seção do relatório descritivo. Vale o art. 24 da LPI: %%
%% tudo o que for necessário para que alguém reproduza o objeto precisa %%
%% estar AQUI. O examinador avalia a suficiência descritiva por este    %%
%% texto (Resolução INPI/PR nº 124/2013, itens 2.13 a 2.16).            %%
%%                                                                     %%
%% Sinais de referência: sempre que citar um elemento do desenho, cite  %%
%% o sinal entre parênteses — "a haste de acionamento (4) desloca o     %%
%% assento (9)". Os mesmos sinais devem identificar o mesmo elemento em %%
%% todas as figuras (art. 39, IV) e ser uniformes em todo o pedido      %%
%% (art. 23, IV).                                                       %%
%%                                                                     %%
%% Unidade do objeto:                                                   %%
%%   invenção — um único conceito inventivo (art. 22 da LPI). Motor novo %%
%%     e sistema de freios novo resolvem problemas técnicos diferentes: %%
%%     são dois pedidos.                                                %%
%%   modelo de utilidade — uma única unidade técnico-funcional e corporal%%
%%     (art. 23 da LPI). Mesa nova e cadeira nova são dois corpos com   %%
%%     funções distintas: são dois pedidos.                             %%
%%                                                                     %%
%% Atenção (art. 27, VII): em modelo de utilidade NÃO cabem trechos do  %%
%% tipo "concretização preferida" ou "a título exemplificativo" — o     %%
%% texto deve definir o objeto requerido e suas variantes, não um       %%
%% princípio que admita formas diversas.                                %%
%%%%%%%%%%%%%%%%%%%%%%%%%%%%%%%%%%%%%%%%%%%%%%%%%%%%%%%%%%%%%%%%%%%%%%%%

\secaoinpi{Descrição \danatureza}

\pnum O dispositivo de acionamento (1) compreende um corpo tubular (2), provido
de bocal de entrada e de bocal de saída, no interior do qual se aloja uma
membrana elástica (3) solidária a uma haste de acionamento (4).

\pnum A haste de acionamento (4) apresenta, em sua porção intermediária, um
ressalto de encosto que atua sobre uma mola de retorno (5) pré-carregada,
apoiada contra a face interna do corpo tubular (2). Na extremidade oposta à
membrana elástica (3), a haste de acionamento (4) recebe um assento de vedação
de dupla face (9), que veda alternadamente o bocal de saída e um orifício de
alívio do corpo tubular (2).

\pnum Um atuador eletromecânico (6) acopla-se à haste de acionamento (4) e
promove o deslocamento desta entre duas posições extremas. A pré-carga da mola
de retorno (5) é dimensionada de modo que, ultrapassado o ponto de inversão
definido pelo ressalto de encosto, a mola de retorno (5) conduza a haste de
acionamento (4) até a posição extrema seguinte e a mantenha nessa posição sem
alimentação do atuador eletromecânico (6).

\pnum Um sensor de pressão (7), montado no bocal de saída do corpo tubular (2),
fornece o sinal de pressão de linha a uma unidade de controle (8), que comanda
o atuador eletromecânico (6) por meio de pulsos de corrente de duração
determinada.

\pnum No estado de abertura, o assento de vedação de dupla face (9) mantém
vedado o orifício de alívio e liberado o bocal de saída, e a água escoa da
entrada para a saída do corpo tubular (2). No estado de fechamento, o assento
de vedação de dupla face (9) mantém vedado o bocal de saída e liberado o
orifício de alívio, o que despressuriza a câmara sob a membrana elástica (3) e
consolida o fechamento.

\pnum A transição entre os dois estados ocorre exclusivamente durante o pulso
de corrente aplicado ao atuador eletromecânico (6), cuja duração é inferior a
duzentos milissegundos, de modo que o consumo de energia por ciclo de
acionamento se reduz à energia do pulso.

\pnum O processo de controle de vazão executado pela unidade de controle (8)
compreende as etapas de: leitura do sinal de pressão fornecido pelo sensor de
pressão (7); comparação do valor lido com uma faixa de operação previamente
definida; e aplicação de um pulso de corrente ao atuador eletromecânico (6)
quando o valor lido permanecer fora da faixa de operação por intervalo superior
a um tempo de confirmação.

\pnum A etapa de comparação emprega uma faixa de operação delimitada por um
limite inferior e um limite superior de pressão, e o tempo de confirmação
impede que oscilações transitórias da linha provoquem acionamentos
desnecessários.
